\section{Discussion}
\label{sec:unit_discussion}
\paragraph{Consistent units in integration and retrieval.}
A standardized property is associated with a canonical unit. Source fields describing the same physical quantity are aligned to this property even when their reported units differ. For example, lengths reported in metres and centimetres share one property with metres as its canonical unit. The source mapping retains the original value and unit together with the matched property, whose unit registry specifies the target unit. Before insertion into the EAV value table, a deterministic conversion maps each measurement to the canonical scale. Query processing applies the same conversion to numerical thresholds after field matching and before SQL compilation. Thus, a condition such as length greater than 120~cm is evaluated as length greater than 1.2~m against values stored in metres. The conversion supports multiplicative scales and affine offsets, such as Celsius to kelvin. Units therefore provide numerical conversion metadata alongside semantic field alignment. The conversion layer is an implementation extension; the reported experiments use the frozen benchmark database.
